\chapter{Carbamazepine Level}

\section*{What Is This Test?}
The carbamazepine test measures the concentration of carbamazepine (CBZ) in the blood. Carbamazepine is a tricyclic anticonvulsant used for epilepsy (focal and generalized tonic-clonic seizures), trigeminal and glossopharyngeal neuralgia, and bipolar disorder (mania). It works primarily by blocking voltage-gated sodium channels in the inactivated state, which stabilizes neuronal membranes and reduces repetitive neuronal firing. Carbamazepine has a narrow therapeutic index: the therapeutic range is 4--12 mcg/mL, and toxicity develops at levels above 12--15 mcg/mL. Therapeutic drug monitoring (TDM) of carbamazepine is standard practice because of its dose-dependent toxicity, autoinduction of its own metabolism, and significant inter-patient variability in pharmacokinetics. Carbamazepine is metabolized by CYP3A4 (and to a lesser extent CYP2C8) to its active metabolite carbamazepine-10,11-epoxide (CBZE), which contributes to both efficacy and toxicity. Carbamazepine induces CYP3A4 and other drug-metabolizing enzymes (autoinduction), so a given dose produces progressively lower levels over the first 3--4 weeks of therapy; dose adjustments are typically needed.

\section*{Alternative Names}
\begin{itemize}
  \item Tegretol Level
  \item Carbatrol Level
  \item Carbamazepine Serum Concentration
  \item Carbamazepine-10,11-Epoxide
\end{itemize}
\section*{Principle}
Carbamazepine is measured by immunoassay (CEDIA, EMIT, or chemiluminescent methods) or by liquid chromatography-tandem mass spectrometry (LC-MS/MS). The specimen is a serum or plasma sample collected in a plain (red-top) or lithium heparin (green-top) tube. Levels are reported in mcg/mL. The therapeutic range is 4--12 mcg/mL. Trough levels (drawn immediately before the next dose) are preferred for monitoring. Carbamazepine exhibits non-linear (saturable) metabolism at higher doses, so small dose increases can produce large and unpredictable increases in serum levels. The free (unbound) fraction is 20--30\%; it is increased in hypoalbuminemia, pregnancy, and uremia, which can produce toxicity at normal total levels. The active metabolite carbamazepine-10,11-epoxide is normally 10--15\% of the parent level but is increased with concurrent valproate or CYP3A4 inhibition, and contributes to toxicity. Carbamazepine induces its own metabolism (autoinduction), so levels fall during the first weeks of therapy and then stabilize; the half-life decreases from 25--65 hours (initial) to 12--17 hours (chronic therapy).

\section*{History}
Carbamazepine was synthesized in 1953 by Walter Schindler at Geigy Pharmaceuticals. It was introduced as an antiepileptic drug in the 1960s and approved for trigeminal neuralgia and epilepsy. Its efficacy for bipolar disorder was established in the 1970s--1980s. Therapeutic drug monitoring of carbamazepine was established early because of its narrow therapeutic index and dose-dependent toxicity \cite{patsalos2008}. The discovery of the strong association between HLA-B*1502 and carbamazepine-induced Stevens-Johnson syndrome (SJS) in Han Chinese patients (Chung and colleagues, 2004) transformed pharmacogenetic screening for carbamazepine \cite{chung2004}. Contemporary guidelines (e.g., NICE 2021, and the guidelines on epilepsy from ILAE) continue to recommend TDM for carbamazepine, particularly during dose titration, pregnancy, and when interacting drugs are added \cite{kwan2011}.

\section*{Why This Test Is Ordered}
TDM every 3--6 months during stable therapy. Monitoring during dose changes or interacting drug changes. Diagnosis of carbamazepine toxicity. Pre-screening for HLA-B*15:02 in at-risk Asian populations. Additional uses: evaluation of non-adherence or suspected subtherapeutic levels (a low level with breakthrough seizures may indicate under-dosing, malabsorption, or rapid metabolism), and monitoring during pregnancy, when carbamazepine levels fall as clearance and volume of distribution increase.

\section*{Primary vs Secondary Test}
The carbamazepine blood level is a primary monitoring test (TDM). It does not diagnose disease; it is used to optimize therapy and prevent toxicity.

\section*{How This Test Is Performed}
A healthcare professional draws blood from a vein. For a trough level, the sample is drawn immediately before the next scheduled dose (typically in the morning, before the first dose of the day). For a peak level (rarely needed), the sample is drawn 2--4 hours after the dose. The sample is collected in a serum separator tube (gold-top), centrifuged, and analyzed by immunoassay or LC-MS/MS. Results are typically available within 1--2 hours for urgent monitoring, or within 24 hours for routine monitoring.

\section*{What Preparation Is Needed}
No fasting is required. The patient should take their carbamazepine as prescribed and note the time of the last dose; the laboratory and provider need the dose time to interpret the result. Document all current medications, especially other anticonvulsants (phenytoin, valproate, phenobarbital, lamotrigine), oral contraceptives, warfarin, and CYP3A4 inhibitors (erythromycin, clarithromycin, ketoconazole, fluconazole, verapamil, diltiazem, grapefruit juice) and inducers (rifampin, phenytoin, phenobarbital, St. John's wort).

\section*{Contraindications and Precautions}
\begin{itemize}
  \item Interpret the result against the time of sampling and the dosing schedule; only a true trough is valid for TDM. Random levels are difficult to interpret because levels fluctuate between doses
  \item Autoinduction lowers carbamazepine levels over the first 3--6 weeks of therapy; repeat the level after dose stabilization, and do not misread this fall as non-adherence
  \item Carbamazepine is a potent inducer of CYP3A4 and other enzymes and lowers levels of oral contraceptives, warfarin, cyclosporine, tacrolimus, and many other drugs; recheck levels after any interacting drug change
  \item Carbamazepine-10,11-epoxide, the active metabolite, contributes to neurotoxicity (ataxia, dizziness, diplopia) but is not measured by standard immunoassays; consider LC-MS/MS measurement if epoxide toxicity is suspected (e.g., with concurrent valproate)
  \item Evaluate for toxicity if nystagmus, ataxia, diplopia, or drowsiness develop. Total levels may be misleading in hypoalbuminemia, pregnancy, and uremia, where the free (active) fraction is increased; check a free carbamazepine level if serum protein is abnormal
  \item HLA-B*15:02 screening is recommended in at-risk Asian populations (Han Chinese, Thai, Malaysian, Filipino, Southern Indian) before initiation \cite{chung2004}, and CYP3A4 inhibitors (macrolides, azole antifungals, verapamil, grapefruit juice) that raise carbamazepine levels require dose adjustment
  \item Monitor CBC, serum sodium, and liver function tests for blood dyscrasias, hyponatremia (SIADH-like effect in 5--30\% of patients), and hepatotoxicity. Carbamazepine is also contraindicated in acute intermittent porphyria and with monoamine oxidase inhibitors (MAOIs)
\end{itemize}
\section*{Risks}
Standard venipuncture risks: mild pain, bruising, or hematoma. The main risk of the test itself is a misleading result (wrong timing relative to the dose), which can lead to inappropriate dose changes.

\section*{What Happens After the Test}
The level is interpreted in the context of the dose, timing, and clinical response. A level in the therapeutic range (4--12 mcg/mL) with good seizure control and no side effects indicates an adequate dose. A level above the therapeutic range with symptoms of toxicity warrants a dose reduction. A level below the therapeutic range with breakthrough seizures warrants a dose increase. Levels are repeated after each dose adjustment (after 2--4 weeks to reach steady state).

\section*{Understanding Results}
Subtherapeutic (<4 mcg/mL): inadequate seizure control or bipolar mood stabilization; autoinduction has lowered levels. Causes include under-dosing, non-adherence, malabsorption, and concurrent enzyme-inducing drugs (phenytoin, phenobarbital, rifampin); dose increase is usually indicated. Low levels may also occur during pregnancy (second and third trimesters) because of increased clearance, volume of distribution, and free fraction; the dose may need to be increased by 50--100\%, with levels monitored each trimester. Low levels with breakthrough seizures warrant gradual dose increases (usually 200 mg/day increments every 1--2 weeks) with level rechecking. Toxic (>12 mcg/mL): nystagmus, ataxia, diplopia, dizziness, sedation; dose reduction is usually indicated. Above 15 mcg/mL, ataxia, diplopia, nystagmus, dysarthria, vomiting, and marked drowsiness occur. Severe toxicity (>20 mcg/mL): coma, seizures, respiratory depression, cardiac conduction abnormalities. In acute overdose, levels above 40 mcg/mL are a medical emergency: carbamazepine is highly protein-bound, so hemodialysis is relatively ineffective; supportive care and multi-dose activated charcoal are the mainstays of treatment, and the level may continue to rise for 12--24 hours after ingestion because the drug's anticholinergic effects delay gastric emptying. With a therapeutic total level and toxic symptoms, consider concurrent epoxide (CBZE) elevation (valproate co-administration, hepatic disease), free-level elevation (hypoalbuminemia, pregnancy), or a drug interaction (verapamil, erythromycin, and grapefruit juice inhibit CYP3A4); measure free carbamazepine or CBZE if available.
\section*{Normal Results}
Therapeutic range: 4--12 mcg/mL (total) for epilepsy and trigeminal neuralgia; a level of 4--8 mcg/mL may be adequate for trigeminal neuralgia. Toxicity is usually seen above 12--15 mcg/mL (symptoms may appear at lower levels in some patients). Free carbamazepine: 1--4 mcg/mL (approximately 20--30\% of the total level; the free fraction is increased in hypoalbuminemia, pregnancy, and uremia). Carbamazepine-10,11-epoxide: 10--15\% of the parent level (increased with valproate co-administration). The trough level is drawn immediately before the next dose; steady state is achieved after 3--5 weeks due to autoinduction.
\section*{Follow-up Tests}
Repeat carbamazepine level 2--4 weeks after any dose change (or more frequently during dose titration, pregnancy, or when interacting drugs are added) to confirm a therapeutic steady-state level. Complete blood count (before therapy and periodically) monitors for leukopenia, thrombocytopenia, agranulocytosis, and aplastic anemia. Serum sodium is checked for hyponatremia (SIADH-like effect in 5--30\% of patients), especially in the elderly. Liver function tests (ALT, AST, GGT) and bilirubin are monitored for hepatotoxicity. HLA-B*1502 screening is performed before starting carbamazepine in at-risk Asian populations (Han Chinese, Thai, Malaysian, Filipino, Southern Indian) because of the high risk of Stevens-Johnson syndrome and toxic epidermal necrolysis. Valproate and other anticonvulsant levels are monitored when co-administered, because carbamazepine induces the metabolism of valproate, lamotrigine, ethosuximide, and phenytoin, lowering their levels.

\section*{References}
\bibliographystyle{unsrtnat}
\bibliography{references}

\clearpage
\section*{Regulatory Status and Authorization Date}
This chapter describes a test category, not a specific commercial product. FDA, CDSCO, EU IVDR, and MHRA approval or registration dates therefore cannot be assigned without the assay manufacturer, product name, instrument, and intended use. For laboratory-developed tests, an individual agency approval date may not exist. See the Preface for the interpretation of regulatory status.

\clearpage
